%% ============================================================
%%  Honours Thesis — First Draft
%%  Topic: Cross-Modal Residual Leakage in Multimodal Privacy
%%         Unlearning: Measurement and Inference-Time Suppression
%%
%%  Status: First complete draft — all results are preliminary.
%%          Claims are bounded by current evidence.
%% ============================================================

\documentclass[12pt, a4paper]{article}

% ── packages ─────────────────────────────────────────────────
\usepackage[utf8]{inputenc}
\usepackage[T1]{fontenc}
\usepackage{lmodern}
\usepackage[margin=2.5cm]{geometry}
\usepackage{setspace}
\onehalfspacing
\usepackage{microtype}
\usepackage{amsmath, amssymb}
\usepackage{booktabs}
\usepackage{multirow}
\usepackage{graphicx}
\usepackage{xcolor}
\usepackage[colorlinks=true, linkcolor=blue, citecolor=blue, urlcolor=blue]{hyperref}
\usepackage{natbib}
\bibliographystyle{plainnat}
\usepackage{parskip}
\usepackage{caption}
\usepackage{subcaption}
\usepackage{algorithm}
\usepackage{algorithmic}
\usepackage{enumitem}

% ── custom commands ───────────────────────────────────────────
\newcommand{\todo}[1]{\textcolor{red}{\textbf{[TODO: #1]}}}
\newcommand{\placeholder}[1]{\textcolor{gray}{\textit{[#1]}}}
\newcommand{\leakgap}{\ensuremath{\Delta_{\text{leak}}}}

% ── metadata ─────────────────────────────────────────────────
\title{%
  \textbf{Cross-Modal Residual Leakage in Multimodal\\
  Privacy Unlearning: Measurement and\\
  Inference-Time Suppression}\\[6pt]
  \large Honours Thesis --- First Draft
}
\author{
  \placeholder{Student Name}\\
  \placeholder{University of Sydney}\\
  \placeholder{School of Computer Science}\\[4pt]
  \small Supervised by \placeholder{Supervisor Name}
}
\date{\today}

% ── document ─────────────────────────────────────────────────
\begin{document}

\maketitle
\thispagestyle{empty}

\begin{abstract}
Machine unlearning has emerged as a practical tool for removing private or sensitive
knowledge from large language models (LLMs) after training.
However, current unlearning methods are predominantly evaluated on text-only behaviour,
leaving open a significant vulnerability: even when knowledge appears to be removed from
a model's text-side responses, it may remain accessible through image-conditioned
inputs in multimodal large language models (MLLMs).
This thesis investigates what we term \emph{cross-modal residual leakage} --- the
phenomenon whereby private factual knowledge survives text-side unlearning and can be
partially recovered when a relevant image is provided at inference time.

Using LLaVA-1.5-7B as the experimental platform and a subset of the UnLOK-VQA
benchmark~\citep{unlok2025}, we first establish a concrete, measurable leakage signal:
providing a semantically relevant image (as opposed to a blank dummy image) raises the
model's correct-answer rate from 0.260 to 0.540 on 50 evaluation samples, yielding a
leakage gap of $+0.280$.
We then identify internal hidden-state channels whose activation patterns are most
sensitive to the presence of a real image on \emph{leak-like} cases---samples where the
image, but not the dummy, enables a correct response.
Finally, we implement and evaluate a lightweight inference-time suppression method that
zeros out the most sensitive channels using a forward hook, requiring no retraining of
any model parameters.

Preliminary results on the leak-like subset show that suppressing the top-5 channels in
layer~31 reduces the recovery rate on that subset from $1.000$ to $0.632$, with
seven positive flips and no negative flips within the leak-like set.
However, no stable improvement is observed on the full 50-sample evaluation set,
indicating that the method has targeted suppression potential but does not yet achieve
broad, reliable leakage reduction.
These results are reported honestly as mixed preliminary evidence.
Limitations, open questions, and planned next steps are discussed.
\end{abstract}

\newpage
\tableofcontents
\newpage

% ──────────────────────────────────────────────────────────────
\section{Introduction}
\label{sec:introduction}
% ──────────────────────────────────────────────────────────────

The widespread deployment of large language models (LLMs) in real-world applications
has raised growing concerns about the privacy of information encoded during training.
Models trained on large web-scale corpora may inadvertently memorise personal facts,
confidential information, or copyrighted material, and such knowledge can be elicited
through targeted prompting even without the model's ``awareness'' that it is doing so.
Machine unlearning offers a principled response to this concern: rather than
retraining a model from scratch, unlearning methods modify or constrain a trained model
to suppress knowledge associated with a designated \emph{forget set}, ideally without
degrading performance on unrelated tasks~\citep{llmunlearning2024}.

Despite notable progress, the unlearning literature has raised an important warning: what
appears to be successful forgetting at the surface level may in fact be obfuscation or
suppression that leaves the underlying knowledge partially intact.
\citet{relearning2024} demonstrate that models which pass standard forgetting evaluations
can have their ``forgotten'' knowledge jogged back to the surface through relatively
mild relearning procedures, suggesting that current evaluation protocols may be
insufficient to certify genuine removal.
This brittle quality is not a mere theoretical concern---it motivates asking whether
seemingly unrelated probing pathways, including alternative input modalities, can
also recover suppressed knowledge.

Multimodal large language models (MLLMs), which process both text and images within a
shared representational framework, present a particularly acute version of this
problem.
When an MLLM is subjected to text-side unlearning, the language components of the model
may be modified to reduce textual recall of certain facts.
However, the visual processing pathway and the cross-modal interaction layers are
typically left untouched.
This creates a potential \emph{backpath}: a semantically relevant image provided at
inference time may reactivate knowledge via visual-to-language projections in layers
that were not targeted during unlearning.
The coupling between modalities, as exemplified by large-scale cross-modal alignment
systems~\citep{imagebind2023}, means that effective unlearning in the multimodal setting
demands more than removing information from one modality in isolation.

Recent work has begun to address this gap.
\citet{multidelete2023} introduce the multimodal machine unlearning problem and propose
early methods for forgetting information that spans both vision and language.
\citet{singleimage2024} focus on the efficient unlearning of visual concepts from MLLMs.
\citet{mmunlearner2025} reformulate the problem specifically for instruction-following
MLLMs, noting the need to distinguish between forgetting visual representations and
preserving unrelated textual capabilities.
Nonetheless, the precise mechanism by which visual inputs can bypass text-side
unlearning---and how such leakage might be detected and mitigated---remains an
underexplored question.

This thesis makes the following contributions:

\begin{enumerate}[label=(\arabic*), leftmargin=*, itemsep=2pt]
  \item We define and operationalise \emph{cross-modal residual leakage} as a measurable
    quantity: the increase in a model's correct-response rate when a semantically
    relevant image is substituted for a blank dummy image, evaluated on samples where
    text-side unlearning has nominally removed the target knowledge.

  \item We implement an automated pipeline that identifies \emph{leak-like cases}
    (instances where the image, but not the dummy, enables a correct response) and uses
    hidden-state forward hooks to compute \emph{channel sensitivity scores}---a
    per-channel measure of how strongly each internal dimension responds differently to
    real versus dummy images on those leak-like cases.

  \item We implement a lightweight \emph{inference-time channel suppression} method that
    scales selected high-sensitivity channels to zero using a registered forward hook on
    one transformer layer, requiring no modification of model weights.

  \item We evaluate this method on a 50-sample subset of the UnLOK-VQA
    benchmark~\citep{unlok2025} and report honest preliminary results: the method shows
    a targeted suppression effect on the leak-like subset but does not yet produce
    stable improvement on the full evaluation set.
\end{enumerate}

We are careful throughout to distinguish between what has been implemented, what is
supported by current preliminary evidence, and what remains as future work.
The aim of this thesis is not to claim a fully validated solution but to establish a
grounded framework for thinking about cross-modal leakage and to provide initial evidence
that internal channel analysis is a meaningful approach to studying it.


% ──────────────────────────────────────────────────────────────
\section{Related Work}
\label{sec:related}
% ──────────────────────────────────────────────────────────────

\subsection{Machine Unlearning in Language Models}

Machine unlearning refers to the problem of modifying a trained model so that it no
longer encodes or expresses knowledge associated with a specified forget
set~\citep{llmunlearning2024}.
In the LLM setting, motivating use cases include the removal of personally identifying
information, the suppression of harmful or copyrighted content, and compliance with
data-deletion regulations.
\citet{llmunlearning2024} provide a comprehensive treatment of the unlearning problem
for LLMs, framing it as a post-training alignment step and evaluating multiple forgetting
strategies on practical tasks.

A key challenge in evaluating unlearning quality is distinguishing genuine removal from
superficial suppression.
\citet{tofu2024} introduce the TOFU benchmark, which evaluates unlearning using
counterfactual and probing questions designed to reveal whether the model has
truly forgotten target information or has merely learned to avoid expressing it in
standard forms.
Their results highlight that simple fine-tuning or gradient ascent approaches can produce
models that appear to forget at the surface while retaining the underlying knowledge.
This observation is reinforced by \citet{relearning2024}, who demonstrate that apparently
unlearned models can be induced to recall their forgotten knowledge through brief benign
relearning procedures.
Together, these findings motivate a more adversarial and probe-oriented view of
unlearning evaluation---one that includes cross-modal probing as an additional axis of
vulnerability.

\subsection{Multimodal and MLLM Unlearning}

The multimodal unlearning problem introduces structural complications beyond those of
text-only settings.
\citet{multidelete2023} are among the first to formalise the problem, noting that
knowledge in multimodal systems exists as associations spanning both visual and
linguistic modalities, and that forgetting must therefore operate across modalities
jointly.
A method that removes textual recall of a fact may leave visual retrieval pathways
largely intact, and vice versa.

More recent work has targeted the instruction-following MLLM specifically.
\citet{singleimage2024} propose an efficient fine-tuning approach for forgetting
visual concepts, using gradient-based updates applied at the level of image-conditioned
generation.
\citet{mmunlearner2025} reformulate the problem in the context of modern instruction-tuned
MLLMs, distinguishing between the goals of forgetting visual patterns and preserving
general textual capabilities.
Their framing is conceptually aligned with the present thesis: both treat
the inter-modal coupling of vision-language models as a primary source of difficulty for
selective forgetting.

In the domain of representation-level vision-language alignment, \citet{cliperase2024}
examine the unlearning of visual-textual associations in CLIP-style models.
Although CLIP is architecturally distinct from an autoregressive MLLM, their finding
that association-level suppression can be achieved without wholesale model modification
provides a useful precedent for approaches that target specific representational
components rather than retraining the full model.

\citet{unlok2025} introduce the UnLOK-VQA benchmark specifically for evaluating
multimodal unlearning in MLLMs.
The benchmark provides knowledge-grounded visual question-answering instances,
each associated with an image identifier and a target answer, and supports evaluation
under attack-defense framing.
The experiments in this thesis use a 50-sample subset of the UnLOK-VQA evaluation data
as the primary evaluation setting.

\subsection{Multimodal Representation Spaces}

A key background intuition for cross-modal leakage is the extent to which vision and
language information becomes entangled within a shared representational space.
\citet{imagebind2023} demonstrate that modern large-scale models can learn a single
embedding space that binds representations from images, text, audio, and other
modalities, such that knowledge associated with one modality is often accessible via
another.
While LLaVA-1.5 is not an ImageBind-style model, the same principle applies at the
level of cross-modal projection layers: the visual token representations produced by
the vision encoder are mapped into the language model's embedding space, and
information flows between these representations during generation.
The CLIP encoder underlying LLaVA's vision component~\citep{clip2021} further
establishes tight semantic alignment between visual and textual representations,
which plausibly facilitates the recovery of textual knowledge through visual prompting.

\subsection{Representation-Level Information Removal}

The present method is conceptually related to, though distinct from, approaches that
suppress unwanted information at the representational level.
\citet{nullspace2020} propose iterative nullspace projection as a way to remove the
influence of protected attributes from text representations, demonstrating that
targeted channel-level manipulation can reduce information leakage without destroying
overall representation quality.
The method in this thesis applies a simpler form of suppression---scaling selected
hidden-state channels by a scalar factor using an inference-time hook---and differs
in motivation (cross-modal leakage rather than demographic attribute removal) and scope
(no iterative retraining is involved).


% ──────────────────────────────────────────────────────────────
\section{Problem Definition and Threat Model}
\label{sec:problem}
% ──────────────────────────────────────────────────────────────

\subsection{Setting}

We consider a multimodal large language model $\mathcal{M}$ that has been subjected to
text-side unlearning with respect to a forget set
$\mathcal{F} = \{(q_i, a_i)\}_{i=1}^{N}$ of factual question-answer pairs.
The unlearning procedure is assumed to have reduced the model's tendency to produce the
target answer $a_i$ in response to the question $q_i$ when presented in a
text-only or dummy-image context.
Concretely, we follow the evaluation framing of \citet{unlok2025} and do not apply
a separate unlearning procedure ourselves; instead, we use the benchmark's data to
construct an \emph{apparent forgetting} scenario and measure the extent to which
image input enables recovery.

\subsection{Cross-Modal Residual Leakage}

\begin{definition}[Cross-Modal Residual Leakage]
\label{def:leakage}
Let $\mathcal{M}$ be an MLLM evaluated on a set of factual VQA instances
$\{(q_i, a_i, v_i)\}_{i=1}^{N}$, where $q_i$ is a question, $a_i$ is the target
answer, and $v_i$ is a semantically relevant image.
Let $v_{\emptyset}$ denote a semantically uninformative dummy image (e.g., a blank
white frame).
Define:
\[
  r_{\text{dummy}} = \frac{1}{N} \sum_{i=1}^{N} \mathbf{1}\bigl[\hat{a}_i^{\emptyset}
  \approx a_i\bigr], \qquad
  r_{\text{true}} = \frac{1}{N} \sum_{i=1}^{N} \mathbf{1}\bigl[\hat{a}_i^{v_i}
  \approx a_i\bigr],
\]
where $\hat{a}_i^{v}$ denotes the model's generated answer when image $v$ is provided.
The \emph{leakage gap} is defined as
\[
  \leakgap = r_{\text{true}} - r_{\text{dummy}}.
\]
A positive $\leakgap$ constitutes evidence of cross-modal residual leakage: the
semantically relevant image enables the model to produce correct responses on cases
where the uninformative baseline fails.
\end{definition}

This definition is designed to be conservative and measurable.
It does not require access to the internals of the unlearning procedure and does not
assume that the model was definitively forced to forget the target knowledge.
Rather, it measures the \emph{marginal information contributed by the visual modality}
as a differential rate.

\subsection{Leak-Like Cases}

A critical subset of instances for our analysis is the \emph{leak-like set}:
\[
  \mathcal{L} = \bigl\{ i : \text{image found} \land
  \mathbf{1}[\hat{a}_i^{v_i} \approx a_i] = 1 \land
  \mathbf{1}[\hat{a}_i^{\emptyset} \approx a_i] = 0 \bigr\}.
\]
These are instances where providing the true image, but not the dummy image, enables a
correct answer.
They represent the most direct observable signal of cross-modal leakage.
Note that $|\mathcal{L}|$ is generally much smaller than $N$; in our experiments,
$|\mathcal{L}| = 18$ out of 50 total samples.

\subsection{Threat Model}

We adopt a passive, query-time threat model.
The adversary is assumed to:

\begin{itemize}[itemsep=2pt]
  \item Have black-box query access to an unlearned MLLM.
  \item Be able to supply images alongside text queries at inference time.
  \item Use semantically relevant images (e.g., images of a person or location
    associated with the forgotten fact) rather than adversarially crafted perturbations.
  \item Evaluate model responses using a soft-match criterion (see
    Section~\ref{sec:setup:eval}) to determine whether the target answer was recovered.
\end{itemize}

This threat model is deliberately mild: the adversary does not perform gradient-based
optimisation, does not have access to model weights, and does not use adversarial images.
We focus on the simplest, most natural form of cross-modal probing because even this
mild setting already reveals a meaningful leakage gap in preliminary experiments.
Stronger adversarial settings---such as adversarial visual prompts or multi-turn
elicitation---are left for future work.

\subsection{Goal of Suppression}

Given the threat model above, the goal of inference-time suppression is to reduce
$\leakgap$ while preserving $r_{\text{dummy}}$ as a proxy for general utility.
Concretely, a successful suppression method should:
\begin{itemize}[itemsep=2pt]
  \item Decrease $r_{\text{true}}^{\text{suppressed}}$ relative to $r_{\text{true}}$,
    ideally approaching $r_{\text{dummy}}$.
  \item Not substantially decrease $r_{\text{dummy}}$, since this would imply
    degradation on non-leaked knowledge.
  \item Show this effect specifically and preferentially on $\mathcal{L}$ rather than
    degrading general outputs indiscriminately.
\end{itemize}
Whether the current method achieves these goals is discussed in
Section~\ref{sec:results}.


% ──────────────────────────────────────────────────────────────
\section{Method Overview}
\label{sec:method}
% ──────────────────────────────────────────────────────────────

The proposed pipeline consists of four stages:
(1) baseline evaluation to measure the leakage gap;
(2) leak-like case filtering;
(3) channel sensitivity analysis via hidden-state forward hooks; and
(4) inference-time selective channel suppression.
All stages operate at inference time: no modification of model weights occurs at any
point.

\begin{figure}[h]
\centering
% \includegraphics[width=0.9\textwidth]{figures/pipeline.pdf}
\placeholder{Pipeline figure: four boxes (Baseline Eval → Leak-Like Filter → Channel
Sensitivity Analysis → Suppression Eval) with arrows. To be added.}
\caption{Overview of the four-stage evaluation pipeline.
No model weights are modified at any stage.}
\label{fig:pipeline}
\end{figure}

\subsection{Stage 1: Baseline Evaluation}

For each sample $(q_i, a_i, v_i)$ in the evaluation set, we run two inference passes
using LLaVA-1.5-7B: one with the dummy image $v_{\emptyset}$ and one with the true
image $v_i$.
Answers are generated using greedy decoding with a maximum of 50 new tokens.
We record whether each generated answer matches the target via a soft-match criterion
(Section~\ref{sec:setup:eval}).
The resulting per-sample records, including match flags for both conditions, are stored
as a JSONL file and used in all downstream stages.

\subsection{Stage 2: Leak-Like Case Filtering}

From the baseline JSONL, we filter for leak-like cases as defined in
Section~\ref{sec:problem}: instances where the image is available, the dummy-image
response does not match the target, and the true-image response does.
These cases form the primary analytical set for channel sensitivity analysis.
The filter is applied programmatically and the resulting subset size is reported in
experiments.

\subsection{Stage 3: Channel Sensitivity Analysis}

For each leak-like case $i \in \mathcal{L}$, we run two additional forward passes
(with $v_{\emptyset}$ and $v_i$) while registering output hooks on selected transformer
layers of the language model decoder.
The hooks capture the full hidden state tensor $\mathbf{h} \in \mathbb{R}^{B \times T \times D}$
at each hooked layer, where $B = 1$ is the batch size, $T$ is the sequence length,
and $D = 4096$ is the hidden dimension for LLaVA-1.5-7B.

The \emph{channel sensitivity score} for channel $c$ at layer $\ell$ is defined as:
\begin{equation}
  s_{\ell, c} = \frac{1}{|\mathcal{L}|} \sum_{i \in \mathcal{L}}
    \frac{1}{T_i} \sum_{t=1}^{T_i}
    \bigl| h^{v_i}_{\ell, t, c} - h^{v_{\emptyset}}_{\ell, t, c} \bigr|,
  \label{eq:sensitivity}
\end{equation}
where $h^v_{\ell, t, c}$ is the value of channel $c$ at token position $t$ in layer
$\ell$ when image $v$ is used.
We compute $s_{\ell, c}$ for all channels and rank them in descending order to obtain
a per-layer top-$k$ list of sensitive channels.
The analysis is run on layers 28 and 31 of the LLaVA-1.5-7B language model backbone.

This analysis is fully inference-time: the model weights are not updated, and the hook
is removed after each forward pass.

\subsection{Stage 4: Inference-Time Channel Suppression}

Given the top-$k$ channel indices $\mathcal{C}_{\ell,k}$ for a target layer $\ell$,
we implement suppression by registering a persistent forward hook during evaluation
that modifies the layer output in-place:
\begin{equation}
  h'_{\ell, :, :, c} \leftarrow \alpha \cdot h_{\ell, :, :, c}
  \quad \forall c \in \mathcal{C}_{\ell, k},
  \label{eq:suppression}
\end{equation}
where $\alpha \in [0, 1]$ is a scalar suppression strength.
At $\alpha = 0.0$, selected channels are zeroed out entirely.
At $\alpha = 1.0$, the hook is a no-op and the model behaves as without suppression.
Intermediate values of $\alpha$ implement partial suppression.

The hook is implemented as a Python context manager using PyTorch's
\texttt{register\_forward\_hook} API, allowing it to be applied cleanly around any
\texttt{model.generate()} call without persistent side effects.
No model parameters are modified, and the suppression is fully reversible.

The method is related in spirit to representation-level information removal approaches
such as nullspace projection~\citep{nullspace2020}, but is simpler and operates at
inference time rather than requiring iterative training.
It differs from fine-tuning-based MLLM unlearning methods~\citep{singleimage2024,
mmunlearner2025} in that it does not update model weights and therefore has no risk of
catastrophic forgetting of unrelated knowledge through gradient updates.

\subsection{Computational Considerations}

All experiments are run on a single GPU (Google Colab, NVIDIA A100 or equivalent)
using 8-bit quantisation via \texttt{bitsandbytes}.
The model is loaded once per evaluation run; the suppression hook adds no meaningful
per-forward-pass overhead.
Channel sensitivity analysis requires two forward passes per leak-like case (dummy and
true image), but does not require generation---only a single full forward pass is needed
to capture hidden states.


% ──────────────────────────────────────────────────────────────
\section{Experimental Setup}
\label{sec:setup}
% ──────────────────────────────────────────────────────────────

\subsection{Model}

All experiments use LLaVA-1.5-7B~\citep{clip2021}, a 7-billion-parameter
multimodal language model consisting of a CLIP ViT-L/14 vision encoder, a two-layer
MLP projection, and a Vicuna-1.5-7B language model backbone (32 transformer layers,
hidden dimension 4096).
The model is loaded from the Hugging Face Hub
(\texttt{llava-hf/llava-1.5-7b-hf}) in 8-bit quantised form.
The model has not been subjected to any unlearning procedure in our experiments;
rather, we use the existing model as a stand-in for a post-unlearning system and
evaluate cross-modal recovery as described below.
All model weights are frozen throughout.

\subsection{Dataset}

We use a 50-sample subset of the UnLOK-VQA evaluation data~\citep{unlok2025}.
UnLOK-VQA provides knowledge-grounded visual question-answering instances, each
consisting of a factual question, a target answer (and associated aliases), and a COCO
image identifier.
Images are sourced from COCO 2017 (train and validation splits).
All 50 images were successfully located; no image-loading errors occurred.

\subsection{Dummy Image}

The dummy image is a $336 \times 336$ blank white RGB frame
($\mathtt{Image.new}(\text{``RGB''}, (336, 336), (255, 255, 255))$).
This is intentionally uninformative: it occupies the image input slot without providing
any semantic content.
The dummy image serves as a no-information baseline to isolate the marginal contribution
of a semantically relevant image to the model's correct-answer rate.

\subsection{Prompt Format}

All inference is performed using the LLaVA chat template:
\begin{verbatim}
USER: <image>
{question}
ASSISTANT:
\end{verbatim}
where \texttt{<image>} is the special image token and \texttt{\{question\}} is the
factual question from the dataset.

\subsection{Evaluation Criterion}
\label{sec:setup:eval}

We evaluate answer correctness using a \emph{soft match} criterion: the generated
answer is considered correct if any of the target answer's aliases (as provided by the
UnLOK-VQA dataset) appears as a normalised substring of the cleaned model output.
Both answer and aliases are lower-cased and stripped of punctuation before comparison.
This criterion is intentionally lenient to account for natural variation in how a model
phrases a correct factual answer.

\subsection{Layer Selection and Suppression Hyperparameters}

Channel sensitivity analysis is performed on layers 28 and 31 of the language model
backbone.
These layers were chosen as representative late-stage decoder layers where cross-modal
integration is expected to be substantial.
In the suppression evaluation, we fix the target layer to layer 31, which shows
stronger channel sensitivity scores than layer 28 in our analysis.

The suppression sweep covers:
\[
  k \in \{1, 3, 5, 10\}, \qquad \alpha \in \{0.0, 0.2, 0.5, 1.0\}.
\]
This yields 16 (layer, $k$, $\alpha$) combinations.
We report results for both the full 50-sample evaluation set and the leak-like subset.

\subsection{Sanity Check}

Before reporting suppression results, we verify that setting $\alpha = 1.0$ (no-op)
reproduces the unsuppressed baseline hit rates on both the full set and the leak-like
subset.
This sanity check confirms that the hook implementation and evaluation path are
functioning correctly.


% ──────────────────────────────────────────────────────────────
\section{Preliminary Results}
\label{sec:results}
% ──────────────────────────────────────────────────────────────

All results reported in this section are preliminary.
The evaluation is conducted on a small sample ($N = 50$) and should be treated as
exploratory evidence rather than definitive conclusions.

\subsection{Baseline Leakage Gap}

Table~\ref{tab:baseline} summarises the baseline evaluation results.
On the full 50-sample set, providing a semantically relevant image raises the model's
soft-match hit rate from 0.260 (dummy image) to 0.540 (true image), yielding a leakage
gap $\leakgap = +0.280$.

\begin{table}[h]
\centering
\caption{Baseline evaluation results ($N = 50$, no suppression).}
\label{tab:baseline}
\begin{tabular}{lcc}
\toprule
\textbf{Condition}           & \textbf{Soft Hit Rate} & \textbf{Leakage Gap $\leakgap$} \\
\midrule
Dummy image ($v_{\emptyset}$) & 0.260 & --- \\
True image ($v_i$)            & 0.540 & $+0.280$ \\
\bottomrule
\end{tabular}
\end{table}

Of the 50 samples, 18 constitute leak-like cases: instances where the true image but not
the dummy enables a correct answer.
These 18 cases are the target of the channel sensitivity analysis and the primary
evaluation set for the suppression experiments.

The leakage gap of $+0.280$ is substantial relative to the base rate and provides
clear initial evidence that semantically relevant images can recover factual information
that is not accessible under the dummy-image condition.
We do not claim that this constitutes a direct measurement of a known unlearning
failure; the baseline model has not been explicitly unlearned.
Rather, this result establishes the measurement framework and shows that the leakage
operationalisation is capable of detecting a real signal.

\subsection{Channel Sensitivity Analysis}

Table~\ref{tab:channels} shows the highest-scoring channels identified for layers 28
and 31 under the channel sensitivity analysis described in
Equation~\eqref{eq:sensitivity}.
The analysis was performed over 18 leak-like cases.

\begin{table}[h]
\centering
\caption{Top-5 sensitive channels per layer (layer 31 scores are stronger).
Full top-20 lists are available in \texttt{channel\_sensitivity.json}.}
\label{tab:channels}
\begin{tabular}{cll}
\toprule
\textbf{Layer} & \textbf{Top-5 Channel Indices} & \textbf{Top Sensitivity Score} \\
\midrule
28 & 1512, 2533, 1076, 1415, 2158 & 15.07 \\
31 & 1512, 3241, 1360, 3556, 2660 & 30.30 \\
\bottomrule
\end{tabular}
\end{table}

Layer 31 exhibits markedly stronger sensitivity than layer 28 (top score 30.30 versus
15.07), suggesting that the cross-modal information relevant to leak-like cases is more
concentrated and discriminative at this layer.
Channel 1512 appears in the top-5 for both layers, which may indicate a consistently
image-sensitive internal dimension across the late decoder.
These observations motivate targeting layer 31 for suppression.

\subsection{Sanity Check: Alpha = 1.0}

Setting $\alpha = 1.0$ (suppression hook active but acting as a no-op) reproduces the
unsuppressed baseline on both the full set ($r_{\text{true}} = 0.540$) and the
leak-like subset ($r_{\text{true}} = 1.000$).
This confirms that the hook registration and evaluation path are implemented correctly
and that any effects observed under other $\alpha$ settings are attributable to the
suppression itself.

\subsection{Suppression Sweep: Full Set}

Table~\ref{tab:sweep_full} reports suppression results on the full 50-sample evaluation
set across the swept hyperparameters.

\begin{table}[h]
\centering
\caption{Suppression sweep on the full evaluation set ($N = 50$, layer 31).
$\delta = r_{\text{supp}} - r_{\text{true}}$.
A negative $\delta$ indicates the suppressed rate is lower than the unsuppressed rate
(desired direction); a positive $\delta$ is undesired.}
\label{tab:sweep_full}
\begin{tabular}{cccccc}
\toprule
$k$ & $\alpha$ & $r_{\text{dummy}}$ & $r_{\text{true}}$ & $r_{\text{supp}}$ & $\delta$ \\
\midrule
 1 & 0.0 & 0.260 & 0.540 & \placeholder{x.xxx} & \placeholder{x.xxx} \\
 1 & 0.2 & 0.260 & 0.540 & \placeholder{x.xxx} & \placeholder{x.xxx} \\
 1 & 0.5 & 0.260 & 0.540 & \placeholder{x.xxx} & \placeholder{x.xxx} \\
 1 & 1.0 & 0.260 & 0.540 & 0.540 & 0.000 \\
\midrule
 3 & 0.0 & 0.260 & 0.540 & \placeholder{x.xxx} & \placeholder{x.xxx} \\
 3 & 0.2 & 0.260 & 0.540 & \placeholder{x.xxx} & \placeholder{x.xxx} \\
 3 & 0.5 & 0.260 & 0.540 & \placeholder{x.xxx} & \placeholder{x.xxx} \\
 3 & 1.0 & 0.260 & 0.540 & 0.540 & 0.000 \\
\midrule
 5 & 0.0 & 0.260 & 0.540 & 0.560 & $+0.020$ \\
 5 & 0.2 & 0.260 & 0.540 & \placeholder{x.xxx} & \placeholder{x.xxx} \\
 5 & 0.5 & 0.260 & 0.540 & \placeholder{x.xxx} & \placeholder{x.xxx} \\
 5 & 1.0 & 0.260 & 0.540 & 0.540 & 0.000 \\
\midrule
10 & 0.0 & 0.260 & 0.540 & \placeholder{x.xxx} & \placeholder{x.xxx} \\
10 & 0.2 & 0.260 & 0.540 & \placeholder{x.xxx} & \placeholder{x.xxx} \\
10 & 0.5 & 0.260 & 0.540 & \placeholder{x.xxx} & \placeholder{x.xxx} \\
10 & 1.0 & 0.260 & 0.540 & 0.540 & 0.000 \\
\bottomrule
\end{tabular}
\end{table}

\textbf{Key observation:} The full-set aggregate does not show a stable improvement
under suppression.
The one confirmed data point ($k=5$, $\alpha=0.0$) shows $r_{\text{supp}} = 0.560$,
meaning the suppressed rate is \emph{slightly higher} than the unsuppressed true-image
rate, in the wrong direction.
The remaining cells are to be filled from the sweep output CSV.
We report the current evidence without favourable selection.

This result indicates that the method is \textbf{not yet a full-set improvement}.
On the broader evaluation set, suppressing selected channels does not reliably decrease
the model's ability to produce correct answers under image prompting.
This may reflect the fact that the sensitive channels carry information relevant to
correct answering on many samples, not only leak-like ones, so zeroing them has mixed
effects across the full distribution.

\subsection{Suppression Sweep: Leak-Like Subset}

Table~\ref{tab:sweep_leak} reports suppression results restricted to the 18 leak-like
cases, which represent the primary target of the method.

\begin{table}[h]
\centering
\caption{Suppression sweep on the leak-like subset ($|\mathcal{L}| = 18$, layer 31).
Good flips: $r_{\text{true}}=1 \to r_{\text{supp}}=0$ (desired).
Bad flips: $r_{\text{true}}=0 \to r_{\text{supp}}=1$ (undesired).
Note: $r_{\text{dummy}}$ and $r_{\text{true}}$ are fixed at 0.000 and 1.000
by definition for this subset.}
\label{tab:sweep_leak}
\begin{tabular}{ccccc}
\toprule
$k$ & $\alpha$ & $r_{\text{supp}}$ & Good Flips & Bad Flips \\
\midrule
 1 & 0.0 & \placeholder{x.xxx} & \placeholder{x} & \placeholder{x} \\
 1 & 0.2 & \placeholder{x.xxx} & \placeholder{x} & \placeholder{x} \\
 1 & 0.5 & \placeholder{x.xxx} & \placeholder{x} & \placeholder{x} \\
 1 & 1.0 & 1.000 & 0 & 0 \\
\midrule
 3 & 0.0 & 0.737 & 5 & 0 \\
 3 & 0.2 & \placeholder{x.xxx} & \placeholder{x} & \placeholder{x} \\
 3 & 0.5 & \placeholder{x.xxx} & \placeholder{x} & \placeholder{x} \\
 3 & 1.0 & 1.000 & 0 & 0 \\
\midrule
 5 & 0.0 & 0.632 & 7 & 0 \\
 5 & 0.2 & \placeholder{x.xxx} & \placeholder{x} & \placeholder{x} \\
 5 & 0.5 & \placeholder{x.xxx} & \placeholder{x} & \placeholder{x} \\
 5 & 1.0 & 1.000 & 0 & 0 \\
\midrule
10 & 0.0 & 0.632 & 7 & 0 \\
10 & 0.2 & \placeholder{x.xxx} & \placeholder{x} & \placeholder{x} \\
10 & 0.5 & \placeholder{x.xxx} & \placeholder{x} & \placeholder{x} \\
10 & 1.0 & 1.000 & 0 & 0 \\
\bottomrule
\end{tabular}
\end{table}

\textbf{Key observations on the leak-like subset:}

\begin{enumerate}[label=(\arabic*), itemsep=4pt]
  \item By construction, the leak-like subset has $r_{\text{true}} = 1.000$ (all 18
    cases produce a correct response under the true image) and
    $r_{\text{dummy}} = 0.000$ (none produce a correct response under the dummy image).

  \item Under $(k=5, \alpha=0.0)$, suppression reduces the recovery rate to $0.632$,
    corresponding to seven cases where the model no longer produces a correct
    answer after channel suppression.
    No bad flips (cases where suppression erroneously introduces a new correct match)
    occur within this subset.

  \item The same seven-flip result is obtained at $(k=10, \alpha=0.0)$, suggesting
    that beyond the top-5 channels, additional channels do not contribute further
    suppression on these cases.

  \item At $(k=3, \alpha=0.0)$, five good flips are observed, indicating that the
    top-3 channels account for a substantial portion of the recoverable signal.

  \item The absence of bad flips within the leak-like subset for the strongest
    suppression settings is an encouraging sign that the method is not simply degrading
    general output quality in a way that incidentally affects these cases; rather,
    it appears to specifically disrupt the image-conditioned recovery pathway.
\end{enumerate}

\textbf{Interpretation:}
These results provide \emph{preliminary evidence} that inference-time channel
suppression can partially disrupt the cross-modal recovery pathway for a non-trivial
fraction of leak-like cases.
The best result (seven out of eighteen good flips, $1.000 \to 0.632$) represents a
$36.8\%$ reduction in recovery rate on the leak-like subset.
However, this evidence is tempered by several important limitations, which we discuss in
Section~\ref{sec:limits}.


% ──────────────────────────────────────────────────────────────
\section{Limitations and Next Steps}
\label{sec:limits}
% ──────────────────────────────────────────────────────────────

\subsection{Current Limitations}

\paragraph{Sample size.}
All experiments are conducted on $N = 50$ samples, of which only 18 are leak-like
cases.
Effect sizes are small in absolute terms (e.g., seven flips out of eighteen).
Results may not generalise to larger evaluation sets, and the observed effects could
partly reflect sampling variability.

\paragraph{No established utility baseline.}
We use the dummy-image hit rate ($r_{\text{dummy}} = 0.260$) as a proxy for utility
preservation, on the assumption that suppression should leave dummy-condition behaviour
largely unchanged.
However, this is not a rigorous utility evaluation.
We do not evaluate suppression on a retain set of unrelated knowledge, and we do not
measure text-only generation quality.
Whether the method preserves general language model capabilities under stronger
suppression settings remains unknown.

\paragraph{Full-set aggregate does not improve.}
On the full 50-sample set, no tested combination of layer, $k$, and $\alpha$ produces
a stable reduction in the overall hit rate under true-image prompting.
The targeted effect on leak-like cases appears alongside neutral or slightly harmful
effects on the broader set.
A method that suppresses leakage only on the leak-like subset while leaving other
image-conditioned responses unchanged would still represent a meaningful contribution,
but this remains to be established more carefully.

\paragraph{Single-layer suppression.}
All reported suppression experiments target layer 31 only.
The leakage pathway may involve multiple layers, and suppressing channels at a single
late layer may insufficiently disrupt earlier cross-modal interactions.
Multi-layer suppression (e.g., layers 28 and 31 jointly) has not yet been evaluated.

\paragraph{No unlearning procedure applied.}
The experiments in this thesis use the pre-unlearning LLaVA-1.5-7B model.
We simulate the unlearning scenario by using the UnLOK-VQA benchmark framing, but do
not apply an explicit unlearning method.
Whether the observed leakage gap would persist after a genuine unlearning procedure---
and whether channel sensitivity analysis on an unlearned model would produce different
top channels---remains an open question.

\paragraph{Sweep results partially pending.}
Tables~\ref{tab:sweep_full} and~\ref{tab:sweep_leak} contain placeholder entries
(marked \placeholder{x.xxx}) for sweep combinations whose results are not yet
confirmed in writing.
These will be filled from the output of \texttt{eval\_with\_suppression.py --sweep}
before the thesis is finalised.

\subsection{Planned Next Steps}

\begin{enumerate}[label=(\arabic*), itemsep=4pt]
  \item \textbf{Complete and fill sweep tables.}
    Run the full $4 \times 4$ sweep for layer 31 and populate
    Tables~\ref{tab:sweep_full} and~\ref{tab:sweep_leak} with all confirmed results.

  \item \textbf{Add a retain-set utility check.}
    Identify a set of general-knowledge VQA samples that are unrelated to the
    forget set and evaluate dummy-image hit rate before and after suppression.
    This would provide a more rigorous utility baseline than $r_{\text{dummy}}$ alone.

  \item \textbf{Compare layer 28 versus layer 31.}
    Run the matched $k \in \{1, 3, 5, 10\}$, $\alpha = 0.0$ sweep targeting layer 28
    and compare leak-like subset results against layer 31 to characterise where the
    cross-modal recovery signal is most concentrated.

  \item \textbf{Evaluate a two-layer suppression variant.}
    If the single-layer results are stable and interpretable, test simultaneous
    suppression at layers 28 and 31 to assess whether additional coverage provides
    further leakage reduction without disproportionate utility cost.

  \item \textbf{Conduct experiments on an explicitly unlearned model.}
    Apply a standard LLM unlearning method (e.g., gradient ascent on the forget set)
    to LLaVA-1.5-7B and re-run the full pipeline on the resulting model.
    This would make the threat model more realistic and directly test the hypothesis
    that cross-modal leakage persists after text-side unlearning.

  \item \textbf{Scale up the evaluation set.}
    Extend the evaluation to at least 200 samples to obtain more reliable estimates
    of the leakage gap and suppression effects.
\end{enumerate}


% ──────────────────────────────────────────────────────────────
\section*{Appendix A: Statements Well Supported by Current Evidence}
\addcontentsline{toc}{section}{Appendix A: Evidence Assessment}
\label{app:evidence}
% ──────────────────────────────────────────────────────────────

The following statements are directly supported by the current experimental results and
can be made with reasonable confidence in the thesis:

\begin{enumerate}[leftmargin=*, itemsep=3pt]
  \item A measurable leakage gap of $+0.280$ exists on the 50-sample evaluation set:
    providing the semantically relevant image raises the hit rate from 0.260 to 0.540.
  \item Eighteen out of fifty samples constitute leak-like cases where the true image,
    but not the dummy image, enables a correct response.
  \item The $\alpha = 1.0$ sanity check confirms the hook implementation is correct
    (reproduces unsuppressed baseline).
  \item On the leak-like subset, suppressing the top-5 channels in layer 31 reduces
    the recovery rate from 1.000 to 0.632 with seven good flips and zero bad flips.
  \item Layer 31 shows higher channel sensitivity scores than layer 28
    (top score 30.30 vs.\ 15.07).
  \item Channel 1512 is among the top-5 sensitive channels in both layers 28 and 31.
\end{enumerate}

% ──────────────────────────────────────────────────────────────
\section*{Appendix B: Statements Needing Stronger Citations}
\addcontentsline{toc}{section}{Appendix B: Citations Needed}
\label{app:citations}
% ──────────────────────────────────────────────────────────────

The following statements currently lack sufficiently specific citation support and
should be strengthened before submission:

\begin{enumerate}[leftmargin=*, itemsep=3pt]
  \item The claim that \emph{``late decoder layers are where cross-modal integration
    is most substantial''} is stated without a mechanistic citation.
    A paper on LLaVA's internal information flow or a probing analysis of late-layer
    behaviour would strengthen this.
  \item The claim that CLIP-based vision encoders facilitate \emph{``tight semantic
    alignment''} that enables factual recovery needs a more specific empirical
    reference, beyond \citet{clip2021}'s general alignment result.
  \item The argument that channel suppression is \emph{``fully reversible and has no
    risk of catastrophic forgetting''} is stated informally and would benefit from a
    citation to the catastrophic forgetting literature for contrast.
  \item The characterisation of the threat model as \emph{``mild and natural''}
    relative to adversarial image attacks would be strengthened by citing at least
    one representative adversarial visual prompt paper.
\end{enumerate}

% ──────────────────────────────────────────────────────────────
\section*{Appendix C: Statements That Should Be Softened}
\addcontentsline{toc}{section}{Appendix C: Softening Needed}
\label{app:softening}
% ──────────────────────────────────────────────────────────────

The following statements in the current draft are either slightly overclaimed or
could be read as stronger than the evidence warrants:

\begin{enumerate}[leftmargin=*, itemsep=3pt]
  \item \emph{``The leakage gap of $+0.280$ is substantial.''} The word
    ``substantial'' is subjective. Prefer: \emph{``a leakage gap of $+0.280$ is
    observed, which is large relative to the dummy-image base rate of 0.260.''} Do not
    imply statistical significance has been established.
  \item \emph{``The absence of bad flips is an encouraging sign that the method is
    not simply degrading general output quality.''} This is too strong given the
    small sample size. Prefer: \emph{``within the leak-like subset, no bad flips
    were observed; however, the sample size is small and this observation should
    be interpreted cautiously.''}
  \item Any reference to the method showing \emph{``targeted suppression potential''}
    should be hedged as \emph{``preliminary evidence of targeted suppression potential
    on a small evaluation subset.''} Do not use language that implies the
    generalisability of the finding.
  \item The contribution statements in the Introduction list ``implement'' and
    ``evaluate'' as completed contributions. The evaluation is preliminary.
    Prefer: \emph{``we implement and provide an initial evaluation of\ldots''}.
\end{enumerate}


% ──────────────────────────────────────────────────────────────
\section*{Appendix D: Supervisor Progress Email Draft}
\addcontentsline{toc}{section}{Appendix D: Supervisor Email Draft}
\label{app:email}
% ──────────────────────────────────────────────────────────────

\begin{quote}
\noindent
\textbf{Subject:} Honours Thesis Progress Update — Cross-Modal Leakage Experiments

\medskip
\noindent
Dear \placeholder{Supervisor name},

I wanted to give you a brief update on where the thesis experiments and writing
currently stand.

On the experimental side, I have completed three of the four pipeline stages.
The baseline evaluation on 50 UnLOK-VQA samples shows a leakage gap of $+0.280$
(dummy hit rate 0.260, true-image hit rate 0.540).
Channel sensitivity analysis on 18 leak-like cases identifies layer 31 as the
primary target, with a top-channel score of 30.30 versus 15.07 for layer 28.
I have implemented inference-time channel suppression and run a preliminary
sweep over $k \in \{1,3,5,10\}$ and $\alpha \in \{0.0,0.2,0.5,1.0\}$ on layer 31.

The headline result at this stage is that suppression does not improve the full-set
aggregate, but on the 18 leak-like cases, suppressing the top-5 channels at
$\alpha=0.0$ reduces the recovery rate from 1.000 to 0.632, with 7 positive flips
and 0 negative flips within that subset.
I am treating this as preliminary mixed evidence and being careful not to overclaim.

On the writing side, I have a first draft covering Introduction, Related Work,
Problem Definition, Method Overview, Experimental Setup, and a Preliminary Results
section.
The results tables have some cells to fill once I confirm the full sweep CSV.

The main gaps I am planning to address next are: (1) a retain-set utility check,
(2) a layer-28 comparison, and (3) scaling up the evaluation set if time allows.

I would very much welcome your feedback on the current draft, particularly on
the framing of the leakage gap operationalisation and whether the contribution
claims are appropriately scoped.
I can share the draft at your convenience.

Best regards,\\
\placeholder{Student name}
\end{quote}


% ──────────────────────────────────────────────────────────────
\section*{Appendix E: Next Experimental Results Needed}
\addcontentsline{toc}{section}{Appendix E: Experimental Roadmap}
\label{app:roadmap}
% ──────────────────────────────────────────────────────────────

The following experimental results are needed to complete or strengthen the thesis:

\begin{itemize}[leftmargin=*, itemsep=4pt]
  \item \textbf{Complete sweep CSV.}
    Fill all placeholder cells in Tables~\ref{tab:sweep_full} and~\ref{tab:sweep_leak}
    by running \texttt{eval\_with\_suppression.py --sweep} and reading the output CSVs.

  \item \textbf{Retain-set utility evaluation.}
    Identify $\approx 50$ general-knowledge VQA samples unrelated to the forget set;
    measure dummy-image hit rate before and after suppression (layer 31, $k=5$,
    $\alpha=0.0$). This directly tests whether suppression preserves utility.

  \item \textbf{Layer 28 suppression comparison.}
    Run the $k \in \{1,3,5,10\}$, $\alpha=0.0$ sweep for layer 28 and compare
    leak-like good flips with layer 31. Determines whether layer targeting matters.

  \item \textbf{Two-layer suppression result (optional).}
    Run one setting (e.g., $k=5$ at both layers 28 and 31, $\alpha=0.0$) to check
    whether joint suppression increases good flips beyond single-layer best.

  \item \textbf{Full sweep for $\alpha \in \{0.2, 0.5\}$ on leak-like subset.}
    Partial-suppression results are currently missing from Table~\ref{tab:sweep_leak}.
    These are needed to characterise the $\alpha$--performance curve.

  \item \textbf{Increased sample size.}
    Re-run the baseline and best suppression setting on $N \geq 200$ samples to
    get more reliable rate estimates and test whether the $+0.280$ gap holds at scale.
\end{itemize}


% ──────────────────────────────────────────────────────────────
%  Bibliography
% ──────────────────────────────────────────────────────────────

\newpage
\bibliography{references}

\end{document}
